\documentclass[12pt,a4paper]{article}

\usepackage[T2A]{fontenc}
\usepackage[utf8]{inputenc}
\usepackage[russian,english]{babel}

\usepackage{amsmath, amssymb, amsthm}
\usepackage{geometry}
\usepackage{array}
\usepackage{booktabs}
\usepackage{listings}
\usepackage{xcolor}
\usepackage{enumitem}

\geometry{left=25mm,right=20mm,top=20mm,bottom=20mm}

\lstdefinestyle{asmstyle}{
    basicstyle=\ttfamily\small,
    keywordstyle=\color{blue},
    commentstyle=\color{gray},
    stringstyle=\color{red!60!black},
    numbers=left,
    numberstyle=\tiny,
    stepnumber=1,
    numbersep=8pt,
    frame=single,
    breaklines=true,
    tabsize=4,
    showstringspaces=false,
    inputencoding=utf8,
    extendedchars=true,
    literate=
        {А}{{\CYRA}}1
        {Б}{{\CYRB}}1
        {В}{{\CYRV}}1
        {Г}{{\CYRG}}1
        {Д}{{\CYRD}}1
        {Е}{{\CYRE}}1
        {Ё}{{\CYRYO}}1
        {Ж}{{\CYRZH}}1
        {З}{{\CYRZ}}1
        {И}{{\CYRI}}1
        {Й}{{\CYRISHRT}}1
        {К}{{\CYRK}}1
        {Л}{{\CYRL}}1
        {М}{{\CYRM}}1
        {Н}{{\CYRN}}1
        {О}{{\CYRO}}1
        {П}{{\CYRP}}1
        {Р}{{\CYRR}}1
        {С}{{\CYRS}}1
        {Т}{{\CYRT}}1
        {У}{{\CYRU}}1
        {Ф}{{\CYRF}}1
        {Х}{{\CYRH}}1
        {Ц}{{\CYRC}}1
        {Ч}{{\CYRCH}}1
        {Ш}{{\CYRSH}}1
        {Щ}{{\CYRSHCH}}1
        {Ъ}{{\CYRHRDSN}}1
        {Ы}{{\CYRERY}}1
        {Ь}{{\CYRSFTSN}}1
        {Э}{{\CYREREV}}1
        {Ю}{{\CYRYU}}1
        {Я}{{\CYRYA}}1
        {а}{{\cyra}}1
        {б}{{\cyrb}}1
        {в}{{\cyrv}}1
        {г}{{\cyrg}}1
        {д}{{\cyrd}}1
        {е}{{\cyre}}1
        {ё}{{\cyryo}}1
        {ж}{{\cyrzh}}1
        {з}{{\cyrz}}1
        {и}{{\cyri}}1
        {й}{{\cyrishrt}}1
        {к}{{\cyrk}}1
        {л}{{\cyrl}}1
        {м}{{\cyrm}}1
        {н}{{\cyrn}}1
        {о}{{\cyro}}1
        {п}{{\cyrp}}1
        {р}{{\cyrr}}1
        {с}{{\cyrs}}1
        {т}{{\cyrt}}1
        {у}{{\cyru}}1
        {ф}{{\cyrf}}1
        {х}{{\cyrh}}1
        {ц}{{\cyrc}}1
        {ч}{{\cyrch}}1
        {ш}{{\cyrsh}}1
        {щ}{{\cyrshch}}1
        {ъ}{{\cyrhrdsn}}1
        {ы}{{\cyrery}}1
        {ь}{{\cyrsftsn}}1
        {э}{{\cyrerev}}1
        {ю}{{\cyryu}}1
        {я}{{\cyrya}}1
}

\lstset{style=asmstyle}

\newtheorem{definition}{Определение}
\newtheorem{theorem}{Теорема}
\newtheorem{remark}{Замечание}

\title{Билет №45 \\ \small Машинно-ориентированные языки, язык ассемблера. Представление машинных команд и констант. Команды транслятору. Их типы, принципы реализации. Макросредства, макровызовы, языки макроопределений, условная макрогенерация, принципы реализации. (СпБГУ)}
\author{Сериков Владислав Александрович}
\date{13 мая 2026}

\begin{document}

\maketitle
\tableofcontents
\newpage

\section{Машинно-ориентированные языки}

\begin{definition}
\textbf{Машинно-ориентированный язык} --- это язык программирования, синтаксис и семантика которого существенно зависят от архитектуры конкретной вычислительной машины: системы команд, регистров, способов адресации, формата машинного слова, организации памяти и соглашений вызова.
\end{definition}

К машинно-ориентированным языкам относят:

\begin{enumerate}
    \item \textbf{машинный язык} --- последовательность двоичных кодов команд и данных, непосредственно исполняемых процессором;
    \item \textbf{язык ассемблера} --- символическое представление машинного языка;
    \item \textbf{автокоды} и специализированные низкоуровневые языки, ориентированные на конкретные ЭВМ;
    \item \textbf{промежуточные низкоуровневые представления}, близкие к машинному коду, например байт-код или LLVM IR в некоторых режимах использования.
\end{enumerate}

Главная особенность машинно-ориентированных языков состоит в том, что программист явно работает с объектами архитектуры: регистрами, ячейками памяти, адресами, флагами процессора, стеком, кодами операций.

\subsection{Уровни представления программы}

Одна и та же программа может рассматриваться на нескольких уровнях:

\begin{center}
\begin{tabular}{lll}
\toprule
Уровень & Пример & Назначение \\
\midrule
Алгоритмический язык & \texttt{x = a + b;} & Удобство для человека \\
Ассемблер & \texttt{add eax, ebx} & Символическая запись машинной команды \\
Машинный код & \texttt{01 D8} & Исполнение процессором \\
Микрооперации & внутренние сигналы CPU & Реализация внутри процессора \\
\bottomrule
\end{tabular}
\end{center}

\section{Язык ассемблера}

\begin{definition}
\textbf{Язык ассемблера} --- машинно-ориентированный язык программирования, в котором машинные команды записываются с помощью мнемоник, а адреса и константы могут задаваться символическими именами.
\end{definition}

Ассемблер является не одним универсальным языком, а семейством языков. Для каждой архитектуры существует собственный набор команд и, как правило, несколько синтаксисов. Например, для x86 распространены Intel-синтаксис и AT\&T-синтаксис.

\subsection{Основные элементы языка ассемблера}

Программа на ассемблере обычно состоит из строк следующих видов:

\begin{enumerate}
    \item \textbf{команды процессора};
    \item \textbf{директивы ассемблера};
    \item \textbf{метки};
    \item \textbf{определения данных};
    \item \textbf{макроопределения};
    \item \textbf{комментарии}.
\end{enumerate}

Общий вид строки ассемблерной программы:

\[
\texttt{[метка:] \quad [операция] \quad [операнды] \quad [; комментарий]}
\]

Пример:

\begin{lstlisting}
start:
    mov eax, 5      ; eax := 5
    add eax, 3      ; eax := eax + 3
\end{lstlisting}

Здесь \texttt{start} --- метка, \texttt{mov} и \texttt{add} --- мнемоники команд, \texttt{eax}, \texttt{5}, \texttt{3} --- операнды.

\subsection{Достоинства и недостатки ассемблера}

\textbf{Достоинства:}

\begin{enumerate}
    \item прямой доступ к аппаратным возможностям;
    \item возможность писать компактный и быстрый код;
    \item полный контроль над регистрами, стеком, памятью;
    \item применимость в системном программировании, загрузчиках, драйверах, встроенных системах;
    \item удобство изучения архитектуры ЭВМ.
\end{enumerate}

\textbf{Недостатки:}

\begin{enumerate}
    \item низкая переносимость;
    \item высокая трудоёмкость разработки;
    \item сложность отладки;
    \item зависимость от конкретной архитектуры и синтаксиса ассемблера;
    \item трудность сопровождения больших программ.
\end{enumerate}

\section{Машинные команды}

\begin{definition}
\textbf{Машинная команда} --- двоично закодированная инструкция процессора, задающая операцию и её операнды.
\end{definition}

Типичная машинная команда содержит:

\begin{enumerate}
    \item \textbf{код операции} --- определяет выполняемое действие;
    \item \textbf{поля регистров} --- указывают используемые регистры;
    \item \textbf{поля способов адресации};
    \item \textbf{смещение или непосредственный операнд};
    \item дополнительные префиксы или модификаторы.
\end{enumerate}

Абстрактно формат команды можно записать так:

\[
\boxed{
\text{opcode}
\mid
\text{mode}
\mid
\text{reg}
\mid
\text{addr}
\mid
\text{immediate}
}
\]

Конкретный формат зависит от архитектуры. В RISC-архитектурах команды часто имеют фиксированную длину, например 32 бита. В CISC-архитектурах, например x86, длина команды может быть переменной.

\subsection{Пример представления команды}

Рассмотрим условный формат команды длиной 16 бит:

\[
\underbrace{0001}_{\text{код операции}}
\underbrace{01}_{\text{регистр 1}}
\underbrace{10}_{\text{регистр 2}}
\underbrace{00000000}_{\text{дополнительное поле}}
\]

Пусть:

\[
0001 = \texttt{ADD}, \qquad
01 = R1, \qquad
10 = R2.
\]

Тогда машинная команда означает:

\[
R1 := R1 + R2.
\]

На ассемблере она может быть записана как:

\begin{lstlisting}
ADD R1, R2
\end{lstlisting}

\section{Представление машинных команд в ассемблере}

Ассемблер заменяет числовые коды операций мнемониками.

\begin{center}
\begin{tabular}{lll}
\toprule
Машинный смысл & Ассемблерная запись & Описание \\
\midrule
Пересылка данных & \texttt{mov eax, ebx} & \texttt{eax := ebx} \\
Сложение & \texttt{add eax, 1} & \texttt{eax := eax + 1} \\
Вычитание & \texttt{sub eax, ebx} & \texttt{eax := eax - ebx} \\
Переход & \texttt{jmp label} & переход к метке \\
Вызов процедуры & \texttt{call f} & вызов подпрограммы \\
Возврат & \texttt{ret} & возврат из подпрограммы \\
\bottomrule
\end{tabular}
\end{center}

\subsection{Операнды команд}

Операнды могут быть:

\begin{enumerate}
    \item \textbf{регистровыми}:
\begin{lstlisting}
mov eax, ebx
\end{lstlisting}

    \item \textbf{непосредственными}:
\begin{lstlisting}
mov eax, 10
\end{lstlisting}

    \item \textbf{операндами в памяти}:
\begin{lstlisting}
mov eax, [x]
\end{lstlisting}

    \item \textbf{адресными выражениями}:
\begin{lstlisting}
mov eax, [ebx + ecx*4 + 8]
\end{lstlisting}
\end{enumerate}

\begin{definition}
\textbf{Непосредственный операнд} --- константа, входящая непосредственно в машинную команду.
\end{definition}

\begin{definition}
\textbf{Адресный операнд} --- операнд, задающий адрес ячейки памяти, из которой нужно прочитать данные или куда нужно записать результат.
\end{definition}

\section{Способы адресации}

\begin{definition}
\textbf{Способ адресации} --- правило вычисления адреса операнда или правило получения самого операнда.
\end{definition}

Основные способы адресации:

\begin{enumerate}
    \item \textbf{непосредственная адресация}:
\begin{lstlisting}
mov eax, 5
\end{lstlisting}
Операндом является сама константа.

    \item \textbf{регистровая адресация}:
\begin{lstlisting}
mov eax, ebx
\end{lstlisting}
Операнд находится в регистре.

    \item \textbf{прямая адресация}:
\begin{lstlisting}
mov eax, [x]
\end{lstlisting}
Адрес задан явно символическим именем.

    \item \textbf{косвенная регистровая адресация}:
\begin{lstlisting}
mov eax, [ebx]
\end{lstlisting}
Адрес находится в регистре \texttt{ebx}.

    \item \textbf{базовая адресация со смещением}:
\begin{lstlisting}
mov eax, [ebx + 4]
\end{lstlisting}

    \item \textbf{индексная адресация}:
\begin{lstlisting}
mov eax, [base + ecx*4]
\end{lstlisting}
Используется, например, для доступа к элементам массива.
\end{enumerate}

\subsection{Пример адресации массива}

Пусть массив целых чисел расположен по адресу \texttt{A}, а каждый элемент занимает 4 байта. Тогда адрес элемента $A[i]$ равен:

\[
\operatorname{addr}(A[i]) = \operatorname{addr}(A) + 4i.
\]

На ассемблере:

\begin{lstlisting}
mov eax, [A + ecx*4] ; eax := A[ecx]
\end{lstlisting}

\section{Представление констант}

\begin{definition}
\textbf{Константа} --- значение, известное во время трансляции программы и не изменяемое в процессе её исполнения.
\end{definition}

В ассемблере константы могут быть:

\begin{enumerate}
    \item числовыми;
    \item символьными;
    \item строковыми;
    \item адресными;
    \item выражениями над другими константами.
\end{enumerate}

\subsection{Системы счисления}

Числовые константы обычно записываются в разных системах счисления.

\begin{center}
\begin{tabular}{lll}
\toprule
Система & Пример & Значение \\
\midrule
Десятичная & \texttt{25} & $25_{10}$ \\
Двоичная & \texttt{11001b} & $25_{10}$ \\
Восьмеричная & \texttt{31o} & $25_{10}$ \\
Шестнадцатеричная & \texttt{19h} или \texttt{0x19} & $25_{10}$ \\
\bottomrule
\end{tabular}
\end{center}

\subsection{Дополнительный код}

Отрицательные целые числа в современных машинах обычно представляются в дополнительном коде.

Для $n$-битного слова число $-x$ представляется как:

\[
2^n - x.
\]

Например, для 8-битного слова:

\[
-5 \equiv 2^8 - 5 = 256 - 5 = 251 = 11111011_2.
\]

Следовательно:

\[
5 = 00000101_2,
\qquad
-5 = 11111011_2.
\]

\subsection{Символьные и строковые константы}

Символьная константа может задаваться кодом символа:

\begin{lstlisting}
letter db 'A'
\end{lstlisting}

Строка --- последовательность байтов:

\begin{lstlisting}
msg db 'Hello', 0
\end{lstlisting}

Здесь нулевой байт часто используется как признак конца строки.

\subsection{Определение данных}

Типичные директивы определения данных:

\begin{center}
\begin{tabular}{lll}
\toprule
Директива & Размер & Пример \\
\midrule
\texttt{db} & 1 байт & \texttt{x db 10} \\
\texttt{dw} & 2 байта & \texttt{y dw 1000} \\
\texttt{dd} & 4 байта & \texttt{z dd 123456} \\
\texttt{dq} & 8 байт & \texttt{q dq 123456789} \\
\bottomrule
\end{tabular}
\end{center}

Пример:

\begin{lstlisting}
section .data
x   db  10
y   dw  1000
z   dd  123456
msg db  'OK', 0
\end{lstlisting}

\section{Команды транслятору}

\begin{definition}
\textbf{Команды транслятору}, или \textbf{директивы ассемблера}, --- конструкции ассемблерного языка, которые не являются командами процессора, а управляют процессом трансляции, размещением данных, генерацией кода и организацией программы.
\end{definition}

Директивы исполняются не процессором, а ассемблером во время трансляции.

\subsection{Отличие команды процессора от директивы}

\begin{center}
\begin{tabular}{lll}
\toprule
Свойство & Команда процессора & Директива ассемблера \\
\midrule
Кто обрабатывает & процессор & ассемблер \\
Когда действует & во время выполнения & во время трансляции \\
Генерирует машинный код & обычно да & не всегда \\
Пример & \texttt{mov eax, ebx} & \texttt{section .data} \\
\bottomrule
\end{tabular}
\end{center}

\section{Типы директив ассемблера}

\subsection{Директивы сегментации и размещения}

Они задают участки программы:

\begin{lstlisting}
section .text
section .data
section .bss
\end{lstlisting}

Обычно:

\begin{enumerate}
    \item \texttt{.text} содержит машинный код;
    \item \texttt{.data} содержит инициализированные данные;
    \item \texttt{.bss} содержит неинициализированные данные.
\end{enumerate}

\subsection{Директивы определения данных}

Пример:

\begin{lstlisting}
x db 10
arr dd 1, 2, 3, 4
\end{lstlisting}

Эти директивы резервируют память и помещают в неё начальные значения.

\subsection{Директивы резервирования памяти}

Пример в синтаксисе NASM:

\begin{lstlisting}
buffer resb 256
array  resd 100
\end{lstlisting}

Здесь \texttt{resb} резервирует байты, а \texttt{resd} --- двойные слова.

\subsection{Директивы присваивания символических имён}

Пример:

\begin{lstlisting}
SIZE equ 100
\end{lstlisting}

После этого имя \texttt{SIZE} может использоваться как константа:

\begin{lstlisting}
mov ecx, SIZE
\end{lstlisting}

\subsection{Директивы внешних и глобальных имён}

Они нужны при раздельной трансляции и компоновке.

\begin{lstlisting}
global main
extern printf
\end{lstlisting}

\texttt{global} делает символ видимым для компоновщика, а \texttt{extern} объявляет символ, определённый в другом модуле.

\subsection{Директивы условной трансляции}

Они позволяют включать или исключать участки текста в зависимости от условий.

Пример:

\begin{lstlisting}
%ifdef DEBUG
    ; отладочный код
%endif
\end{lstlisting}

\subsection{Директивы включения файлов}

Пример:

\begin{lstlisting}
%include "macros.inc"
\end{lstlisting}

Эта директива вставляет содержимое другого файла в текущий исходный текст.

\section{Принципы реализации ассемблера}

Ассемблер должен преобразовать исходный текст в объектный модуль. Главные задачи:

\begin{enumerate}
    \item лексический и синтаксический анализ;
    \item построение таблицы символов;
    \item вычисление адресов;
    \item преобразование мнемоник в коды операций;
    \item вычисление выражений;
    \item генерация машинного кода;
    \item формирование объектного файла;
    \item генерация информации для компоновщика.
\end{enumerate}

\subsection{Проблема прямых ссылок}

В ассемблере возможны ссылки на метки, которые ещё не определены.

Пример:

\begin{lstlisting}
    jmp finish
    mov eax, 1
finish:
    ret
\end{lstlisting}

При обработке команды \texttt{jmp finish} ассемблер ещё не знает адрес метки \texttt{finish}. Поэтому часто используется двухпроходная трансляция.

\section{Двухпроходный ассемблер}

\begin{definition}
\textbf{Двухпроходный ассемблер} --- ассемблер, который просматривает исходную программу дважды: на первом проходе строит таблицу символов и вычисляет адреса, на втором генерирует машинный код.
\end{definition}

\subsection{Алгоритм двухпроходной трансляции}

\textbf{Вход:} исходный текст программы на ассемблере.

\textbf{Выход:} объектный модуль.

\subsubsection*{Первый проход}

\begin{enumerate}
    \item Установить счётчик адреса $LC := 0$.
    \item Читать исходный текст построчно.
    \item Если строка содержит метку, занести пару
    \[
    (\text{имя метки}, LC)
    \]
    в таблицу символов.
    \item Если строка содержит команду, увеличить $LC$ на длину машинной команды.
    \item Если строка содержит директиву определения данных, увеличить $LC$ на размер данных.
    \item Если строка содержит директиву резервирования памяти, увеличить $LC$ на резервируемый размер.
    \item Обработать директивы, влияющие на размещение.
    \item Продолжать до конца файла.
\end{enumerate}

\subsubsection*{Второй проход}

\begin{enumerate}
    \item Снова установить $LC := 0$.
    \item Читать исходный текст построчно.
    \item Для каждой команды определить код операции.
    \item Вычислить операнды, используя таблицу символов.
    \item Сформировать машинное представление команды.
    \item Для директив данных сгенерировать байты данных.
    \item Для внешних ссылок сформировать записи релокации.
    \item Записать результат в объектный модуль.
\end{enumerate}

\subsection{Минимальный пример двухпроходной трансляции}

Исходная программа:

\begin{lstlisting}
start:
    mov eax, 1
    jmp finish
    mov eax, 2
finish:
    ret
\end{lstlisting}

Пусть условные длины команд равны:

\[
|\texttt{mov eax, imm}| = 5,
\qquad
|\texttt{jmp label}| = 5,
\qquad
|\texttt{ret}| = 1.
\]

Первый проход:

\begin{center}
\begin{tabular}{llll}
\toprule
Строка & Действие & $LC$ до & $LC$ после \\
\midrule
\texttt{start:} & метка \texttt{start} = 0 & 0 & 0 \\
\texttt{mov eax, 1} & команда длины 5 & 0 & 5 \\
\texttt{jmp finish} & команда длины 5 & 5 & 10 \\
\texttt{mov eax, 2} & команда длины 5 & 10 & 15 \\
\texttt{finish:} & метка \texttt{finish} = 15 & 15 & 15 \\
\texttt{ret} & команда длины 1 & 15 & 16 \\
\bottomrule
\end{tabular}
\end{center}

Таблица символов:

\[
\begin{array}{c|c}
\text{Символ} & \text{Адрес} \\
\hline
\texttt{start} & 0 \\
\texttt{finish} & 15
\end{array}
\]

На втором проходе команда \texttt{jmp finish} может быть закодирована, поскольку адрес \texttt{finish} уже известен.

\section{Таблица символов}

\begin{definition}
\textbf{Таблица символов} --- структура данных транслятора, сопоставляющая именам их свойства: адрес, тип, область видимости, секцию, размер, признак внешнего или глобального символа.
\end{definition}

Типичная запись таблицы символов:

\[
\langle
\text{name},
\text{value},
\text{section},
\text{type},
\text{binding}
\rangle.
\]

Например:

\begin{center}
\begin{tabular}{llll}
\toprule
Имя & Значение & Секция & Тип \\
\midrule
\texttt{main} & 0 & \texttt{.text} & global \\
\texttt{x} & 0 & \texttt{.data} & local \\
\texttt{printf} & unknown & extern & external \\
\bottomrule
\end{tabular}
\end{center}

\section{Релокация и компоновка}

\begin{definition}
\textbf{Релокация} --- процесс корректировки адресов в объектном коде после определения фактических адресов размещения секций и внешних символов.
\end{definition}

Раздельная трансляция позволяет собирать программу из нескольких объектных модулей. Ассемблер не всегда знает окончательные адреса всех объектов, поэтому он оставляет для компоновщика специальные записи релокации.

Пример:

\begin{lstlisting}
extern f
call f
\end{lstlisting}

Ассемблер знает, что нужно вызвать функцию \texttt{f}, но не знает её адрес. В объектный файл помещается команда вызова и запись релокации, сообщающая компоновщику, что адрес необходимо исправить.

\section{Макросредства}

\begin{definition}
\textbf{Макросредства} --- средства языка ассемблера или препроцессора, позволяющие задавать параметризованные текстовые шаблоны и порождать по ним фрагменты программы.
\end{definition}

Макросредства позволяют:

\begin{enumerate}
    \item сокращать повторяющийся код;
    \item повышать читаемость программы;
    \item описывать псевдокоманды;
    \item реализовывать условную генерацию кода;
    \item параметризовать фрагменты программы.
\end{enumerate}

\section{Макроопределения и макровызовы}

\begin{definition}
\textbf{Макроопределение} --- описание макроса, содержащее имя, список формальных параметров и тело макроса.
\end{definition}

\begin{definition}
\textbf{Макровызов} --- использование имени макроса с фактическими аргументами, приводящее к подстановке тела макроса с заменой формальных параметров фактическими.
\end{definition}

Пример в стиле NASM:

\begin{lstlisting}
%macro clear 1
    xor %1, %1
%endmacro

clear eax
clear ebx
\end{lstlisting}

После макрорасширения получится:

\begin{lstlisting}
xor eax, eax
xor ebx, ebx
\end{lstlisting}

\subsection{Формальные и фактические параметры}

В макроопределении:

\begin{lstlisting}
%macro push2 2
    push %1
    push %2
%endmacro
\end{lstlisting}

\texttt{\%1} и \texttt{\%2} --- формальные параметры.

В макровызове:

\begin{lstlisting}
push2 eax, ebx
\end{lstlisting}

\texttt{eax} и \texttt{ebx} --- фактические параметры.

Результат раскрытия:

\begin{lstlisting}
push eax
push ebx
\end{lstlisting}

\section{Язык макроопределений}

\begin{definition}
\textbf{Язык макроопределений} --- подъязык ассемблера или препроцессора, предназначенный для описания макросов, параметров, локальных меток, условной генерации и повторения фрагментов текста.
\end{definition}

Обычно язык макроопределений включает:

\begin{enumerate}
    \item средства объявления макросов;
    \item позиционные параметры;
    \item именованные параметры;
    \item параметры по умолчанию;
    \item переменное число аргументов;
    \item локальные метки;
    \item условные конструкции;
    \item циклы генерации;
    \item операции над строками и числами.
\end{enumerate}

\subsection{Макрос с локальной меткой}

При многократном раскрытии макроса возникает проблема повторного определения меток.

Плохой пример:

\begin{lstlisting}
%macro loop_once 0
label:
    nop
    jmp label
%endmacro

loop_once
loop_once
\end{lstlisting}

После раскрытия метка \texttt{label} появится дважды.

Поэтому используются локальные или уникализируемые метки:

\begin{lstlisting}
%macro loop_once 0
%%label:
    nop
    jmp %%label
%endmacro
\end{lstlisting}

Метка \texttt{\%\%label} при каждом раскрытии получает уникальное внутреннее имя.

\section{Условная макрогенерация}

\begin{definition}
\textbf{Условная макрогенерация} --- генерация различных фрагментов текста программы в зависимости от условий, вычисляемых во время макрообработки или трансляции.
\end{definition}

Пример:

\begin{lstlisting}
%define DEBUG 1

%if DEBUG
    ; отладочный код
%endif
\end{lstlisting}

Макрос с условной генерацией:

\begin{lstlisting}
%macro save_regs 1
%if %1 = 32
    push eax
    push ebx
%elif %1 = 64
    push rax
    push rbx
%endif
%endmacro
\end{lstlisting}

Вызов:

\begin{lstlisting}
save_regs 64
\end{lstlisting}

Результат:

\begin{lstlisting}
push rax
push rbx
\end{lstlisting}

\section{Принципы реализации макропроцессора}

\begin{definition}
\textbf{Макропроцессор} --- компонент транслятора, выполняющий обработку макроопределений и макровызовов, обычно до основного синтаксического анализа и генерации машинного кода.
\end{definition}

Макропроцессор может быть:

\begin{enumerate}
    \item самостоятельной программой;
    \item частью ассемблера;
    \item частью компилятора;
    \item препроцессором общего назначения.
\end{enumerate}

\subsection{Основные таблицы макропроцессора}

При реализации макропроцессора часто используются:

\begin{enumerate}
    \item \textbf{таблица имён макросов};
    \item \textbf{таблица макроопределений};
    \item \textbf{таблица формальных параметров};
    \item \textbf{стек раскрытия макросов};
    \item \textbf{таблица локальных меток};
    \item \textbf{таблица условной трансляции}.
\end{enumerate}

\subsection{Алгоритм обработки макросов}

\textbf{Вход:} исходная программа с макроопределениями и макровызовами.

\textbf{Выход:} программа после раскрытия макросов.

\begin{enumerate}
    \item Читать входной текст построчно.
    \item Если найдена директива начала макроопределения, считать заголовок макроса.
    \item Занести имя макроса и список параметров в таблицу макросов.
    \item Считать тело макроса до директивы конца макроопределения.
    \item Сохранить тело макроса во внутреннем представлении.
    \item Если найдена обычная строка, проверить, является ли первая операция именем макроса.
    \item Если это не макровызов, передать строку в выходной поток.
    \item Если это макровызов, сопоставить фактические аргументы формальным параметрам.
    \item Скопировать тело макроса в выходной поток, заменяя формальные параметры фактическими.
    \item При необходимости выполнить условные конструкции и циклы макрогенерации.
    \item Уникализировать локальные метки.
    \item Если при раскрытии возник новый макровызов, обработать его рекурсивно или с помощью стека раскрытия.
\end{enumerate}

\subsection{Минимальный пример работы макропроцессора}

Исходный текст:

\begin{lstlisting}
%macro inc2 1
    inc %1
    inc %1
%endmacro

inc2 eax
\end{lstlisting}

Таблица макросов:

\begin{center}
\begin{tabular}{lll}
\toprule
Имя & Число параметров & Тело \\
\midrule
\texttt{inc2} & 1 & \texttt{inc \%1; inc \%1} \\
\bottomrule
\end{tabular}
\end{center}

Макровызов:

\[
\texttt{inc2 eax}
\]

Сопоставление параметров:

\[
\%1 \mapsto \texttt{eax}.
\]

Результат раскрытия:

\begin{lstlisting}
inc eax
inc eax
\end{lstlisting}

\section{Рекурсивные макросы}

Некоторые макропроцессоры допускают рекурсивные макросы.

Пример идеи:

\begin{lstlisting}
%macro repeat_nop 1
%if %1 > 0
    nop
    repeat_nop %1 - 1
%endif
%endmacro
\end{lstlisting}

Такой макрос должен раскрываться с ограничением глубины рекурсии, иначе возможна бесконечная макрогенерация.

\section{Макросы и подпрограммы}

Макросы и подпрограммы решают похожую задачу повторного использования кода, но делают это по-разному.

\begin{center}
\begin{tabular}{lll}
\toprule
Свойство & Макрос & Подпрограмма \\
\midrule
Когда обрабатывается & при трансляции & при выполнении \\
Механизм & текстовая подстановка & переход и возврат \\
Расход памяти на код & может увеличиваться & обычно меньше \\
Расход времени & нет вызова & есть вызов и возврат \\
Параметры & подстановка текста & регистры, стек, память \\
\bottomrule
\end{tabular}
\end{center}

Макросы выгодны для коротких часто используемых фрагментов, где накладные расходы вызова подпрограммы существенны.

\section{Псевдокоманды}

\begin{definition}
\textbf{Псевдокоманда} --- команда, отсутствующая в системе команд процессора, но поддерживаемая ассемблером и раскрываемая в одну или несколько реальных машинных команд.
\end{definition}

Например, если процессор не имеет команды очистки регистра, ассемблер может поддерживать псевдокоманду:

\begin{lstlisting}
clear eax
\end{lstlisting}

которая раскрывается в:

\begin{lstlisting}
xor eax, eax
\end{lstlisting}

Псевдокоманды можно рассматривать как частный случай макросредств или синтаксического сахара.

\section{Формальное описание трансляции ассемблерной программы}

Пусть исходная программа состоит из строк:

\[
P = (l_1, l_2, \ldots, l_n).
\]

Ассемблер строит объектный модуль:

\[
O = A(P),
\]

где $A$ --- функция трансляции.

Трансляция включает два основных отображения:

\[
M: \text{исходный текст} \to \text{текст без макросов},
\]

\[
G: \text{текст без макросов} \to \text{объектный код}.
\]

Тогда:

\[
A = G \circ M.
\]

То есть сначала выполняется макрообработка, затем генерация объектного кода.

\section{Теорема о корректности двухпроходной трансляции}

\begin{theorem}
Пусть ассемблерная программа содержит только метки, директивы размещения, директивы определения данных и команды фиксированной длины или команды, длина которых однозначно определяется без знания значений ещё не определённых меток. Тогда двухпроходный ассемблер, который на первом проходе строит таблицу символов, а на втором генерирует код с использованием этой таблицы, корректно разрешает все внутренние ссылки на метки, определённые в программе.
\end{theorem}

\begin{proof}
Рассмотрим программу как последовательность строк
\[
l_1, l_2, \ldots, l_n.
\]

Пусть счётчик адреса $LC$ в начале первой строки равен начальному адресу секции. По условию длина каждой команды и каждого блока данных может быть определена на первом проходе без знания значений ещё не определённых меток.

На первом проходе ассемблер просматривает строки последовательно. Перед обработкой строки $l_i$ значение $LC$ равно адресу, по которому будет размещён код или данные, порождаемые этой строкой. Если строка содержит метку $s$, то ассемблер заносит в таблицу символов пару
\[
s \mapsto LC.
\]

Поскольку $LC$ в этот момент совпадает с адресом размещения соответствующей строки, в таблицу символов заносится правильный адрес метки.

После этого ассемблер увеличивает $LC$ на размер кода или данных, порождённых строкой. Так как размер каждой строки определяется корректно, по индукции по номеру строки получаем, что перед каждой строкой $l_i$ значение $LC$ равно её фактическому адресу размещения.

Следовательно, после завершения первого прохода таблица символов содержит правильные адреса всех меток, определённых в программе.

На втором проходе ассемблер снова просматривает программу. При встрече ссылки на метку $s$ он обращается к таблице символов. Если метка определена в программе, то по доказанному выше таблица содержит её правильный адрес. Поэтому ассемблер подставляет в машинную команду правильное абсолютное значение адреса или правильное относительное смещение, вычисляемое из правильного адреса метки и текущего адреса команды.

Следовательно, все внутренние ссылки на определённые метки разрешаются корректно. Теорема доказана.
\end{proof}

\section{Итоговая схема обработки ассемблерной программы}

Полная обработка может быть представлена следующей схемой:

\[
\boxed{
\text{исходный текст}
\to
\text{макрообработка}
\to
\text{ассемблирование}
\to
\text{объектный файл}
\to
\text{компоновка}
\to
\text{исполняемый файл}
}
\]

Этапы:

\begin{enumerate}
    \item Макропроцессор раскрывает макросы.
    \item Ассемблер преобразует мнемоники в машинные команды.
    \item Формируется объектный файл с секциями, символами и релокациями.
    \item Компоновщик объединяет объектные файлы.
    \item Загрузчик размещает исполняемый файл в памяти.
    \item Процессор исполняет машинные команды.
\end{enumerate}

\section{Минимальный пример программы}

Пример на x86-64 NASM для Linux:

\begin{lstlisting}
section .data
msg db "Hello", 10
len equ $ - msg

section .text
global _start

_start:
    mov rax, 1      ; номер системного вызова write
    mov rdi, 1      ; файловый дескриптор stdout
    mov rsi, msg    ; адрес строки
    mov rdx, len    ; длина строки
    syscall

    mov rax, 60     ; номер системного вызова exit
    xor rdi, rdi    ; код возврата 0
    syscall
\end{lstlisting}

Здесь:

\begin{enumerate}
    \item \texttt{section .data} задаёт секцию данных;
    \item \texttt{msg db "Hello", 10} определяет строку;
    \item \texttt{len equ \$ - msg} вычисляет длину строки;
    \item \texttt{global \_start} экспортирует точку входа;
    \item \texttt{mov}, \texttt{xor}, \texttt{syscall} являются машинными командами;
    \item \texttt{section}, \texttt{global}, \texttt{equ}, \texttt{db} являются директивами ассемблера.
\end{enumerate}

\section{Краткие выводы}

\begin{enumerate}
    \item Машинно-ориентированные языки тесно связаны с архитектурой ЭВМ.
    \item Ассемблер является символическим представлением машинного языка.
    \item Машинная команда содержит код операции и информацию об операндах.
    \item Константы в ассемблере могут задаваться в разных системах счисления и использоваться как непосредственные операнды или данные.
    \item Директивы ассемблера управляют трансляцией, но не являются командами процессора.
    \item Для разрешения прямых ссылок часто используется двухпроходная трансляция.
    \item Макросредства позволяют порождать фрагменты программы на этапе трансляции.
    \item Макросы отличаются от подпрограмм тем, что раскрываются до выполнения программы.
    \item Условная макрогенерация позволяет создавать разные варианты кода из одного исходного текста.
\end{enumerate}

\end{document}
